% Partie : logique intuitionniste
% Chapitre : propositions comme types
% Section : normalisation

\documentclass[../../../include/open-logic-section]{subfiles}

\begin{document}

\olfileid[fr]{int}{pty}{nor}

\olsection{Normalisation}

Nous allons montrer que toute dérivation bien formée
peut être réduite à une dérivation normale : c'est la
\emph{propriété de normalisation}. Par la correspondance
entre propositions et types, il suffit d'étudier les
termes de preuve bien typés. Il faut distinguer deux
énoncés : l'existence d'une suite de réductions qui
atteint une forme normale et la terminaison de
\emph{toutes} les suites de réductions.

Commençons par des mesures de complexité. La
\emph{longueur}~$\len{!A}$ d'une formule est définie par
\begin{align*}
  \len{p} &= 0\\
  \len{\lfalse} &= 0\\
  \len{!A \land !B} &= \len{!A} + \len{!B} + 1\\
  \len{!A \lor !B} &= \len{!A} + \len{!B} + 1 \\
  \len{!A \lif !B} &= \len{!A} + \len{!B} + 1.
\end{align*}
Ici $p$ est atomique. Pour chacun des trois redex
principaux, le \emph{rang de coupure} est la longueur du
type composé immédiatement introduit puis éliminé :
\begin{align*}
  \cutrank{(\lambd[\typeof{x}{!A}][\typeof{N}{!B}])Q} &=
  \len{!A} + \len{!B} + 1 \\
  \cutrank{\ande{i}{\andi{\typeof{M}{!A}}{\typeof{N}{!B}}}} &=
  \len{!A} + \len{!B} + 1 \\
  \cutrank{\ore{\ori{i}{!A_{3-i}}{\typeof{M}{!A_i}}}
    {\typeof{x_1}{!A_1}}{\typeof{N_1}{!C}}
    {\typeof{x_2}{!A_2}}{\typeof{N_2}{!C}}} &=
  \len{!A_1} + \len{!A_2} + 1.
\end{align*}
La notation d'injection indique le type de l'autre
sommande. Les conversions de permutation ne sont pas
des redex principaux et ne reçoivent pas ce rang.
Posons, pour la famille finie des redex principaux de~$M$,
\[
  \maxrank{M} =
  \max\bigl(\{0\}\cup
    \{\cutrank{N}\mid N\text{ est un sous-terme redex principal de }M\}
  \bigr).
\]
Ainsi, si $M$ n'a aucun redex principal, son rang maximal
vaut~$0$ ; cela n'exclut pas une éventuelle conversion
de permutation.

\begin{lem}\ollabel{lem:subst}
  Supposons que $M$ ne soit pas la seule variable~$x$,
  que $x$ ait le type~$!A$ et que
  $\Subst{M}{\typeof{N}{!A}}{\typeof{x}{!A}}$
  soit un redex principal \emph{à la racine}.
  Ou bien $M$ était déjà un redex principal à la racine,
  de même rang, ou bien l'une des situations suivantes
  crée ce redex :
  \begin{enumerate}
  \item $M$ est de la forme $\ande{i}{x}$ et $N$
    est de la forme $\andi{P_1}{P_2}$ ;
  \item $M$ est de la forme
    $\ore{x}{x_1}{P_1}{x_2}{P_2}$ et $N$
    est de la forme $\ori{i}{!A_{3-i}}{Q}$ ;
  \item $M$ est de la forme $x Q$ et $N$
    est de la forme $\lambd[y][P]$.
  \end{enumerate}
  Dans les trois derniers cas,
  $\cutrank{\Subst{M}{N}{x}}=\len{!A}$.
\end{lem}
\begin{proof}
  Examinons le constructeur extérieur de~$M$.
  Une introduction ne devient pas, par substitution,
  une élimination à la racine. Si la racine de~$M$
  formait déjà un détour introduction--élimination,
  la substitution ne change pas le type de la formule
  intermédiaire, donc ne change pas son rang.
  Sinon, pour que l'élimination extérieure devienne
  un redex principal, sa prémisse majeure doit être
  la variable~$x$, que~$N$ remplace par l'introduction
  correspondante. Les seules possibilités sont
  respectivement la projection, l'analyse par cas
  et l'application énumérées ci-dessus. Dans chaque cas,
  la formule introduite puis éliminée est précisément
  le type~$!A$ de~$x$.
\end{proof}

\begin{lem}
  Soit $R \redone R'$ une contraction principale à la
  racine. Si le rang de chaque redex principal proprement
  contenu dans~$R$ est strictement inférieur à
  $\cutrank{R}$, alors
  $\maxrank{R'}<\cutrank{R}$.
\end{lem}
\begin{proof}
  Si $R$ est la projection d'un couple, $R'$ est une
  composante du couple ; tous ses redex étaient déjà
  des sous-termes propres de~$R$.
  Pour une application de la forme
  $(\lambd[\typeof{x}{!A}][N])\typeof{Q}{!A}$,
  le réduit est $\Subst{N}{Q}{x}$.
  Un redex qui s'y trouve provient soit d'un redex
  propre de~$R$, éventuellement recopié lors de la
  substitution, soit de la rencontre d'une occurrence
  de~$x$ avec une introduction fournie par~$Q$.
  Dans ce dernier cas, le lemme précédent donne le
  rang~$\len{!A}$, strictement inférieur à la longueur
  de~$!A\lif !B$, type de l'abstraction
  immédiatement éliminée.
  Si $R$ est l'analyse d'une injection de type
  $!A_1\lor !A_2$, le réduit est la substitution dans
  la branche choisie. Les redex nouvellement créés
  ont le rang~$\len{!A_i}$, strictement inférieur à
  $\len{!A_1\lor !A_2}$ ; les redex déjà présents dans
  le terme injecté et les branches satisfont
  l'hypothèse du lemme.
\end{proof}

Ces deux lemmes sont \emph{locaux} : ils examinent le
réduit d'un redex isolé. Ils ne fournissent pas une
mesure décroissante pour un terme entier. Par exemple,
réduire une analyse par cas à l'intérieur d'une projection
peut produire un nouveau redex de projection dont le
rang est supérieur à celui de l'analyse réduite.
La récurrence sur le rang maximal et le nombre de redex
de ce rang demanderait de maîtriser aussi les redex
créés au-dessus du sous-terme réduit. Le résultat global
s'obtient par le théorème suivant.

\begin{thm}
  Tout terme de preuve bien typé de ce calcul est
  fortement normalisant pour les contractions principales
  et la permutation d'une projection dans une analyse
  par cas. Toute suite de réductions est finie.
  Par conséquent, toute dérivation correspondante
  se réduit à une dérivation normale.
\end{thm}
\begin{proof}
  Les types et les constructeurs considérés ici sont
  ceux du lambda-calcul simplement typé avec implication,
  produit, somme et type vide. Le théorème de
  \href{https://arxiv.org/abs/0804.2535}{Wojdyga,
  \emph{Short proofs of strong normalization},
  théorème~14, arXiv:0804.2535v1} s'applique à ce calcul.
  Ses définitions~3 et~4 contiennent les contractions
  principales et les deux instances de la permutation
  d'une projection dans une analyse par cas. Une suite
  infinie pour nos règles serait donc une suite infinie
  dans le calcul de ce théorème, ce qu'il exclut.
  Les conversions des dérivations sont reflétées par
%
%
  celles des termes de preuve.
\end{proof}

En particulier, chaque terme bien typé
\emph{possède} une forme normale. Si l'on interprète
la réduction vers cette forme comme l'exécution d'un
calcul, la normalisation forte dit davantage :
quel que soit l'ordre des contractions autorisées,
l'exécution termine. La conservation des types assure
que le résultat garde le type annoncé. Ainsi, si
$N$ a le type $!A \lif !B$ et $M$ le type~$!A$,
toute réduction complète de l'application $(NM)$
se termine en un terme $N'$ de type~$!B$.
Dans cette analogie, « valeur » désigne une forme
normale ; pour un terme ouvert, cette forme peut
encore contenir des variables libres.

La terminaison et la conservation des types ne suffisent
pas à garantir que deux ordres de réduction donnent
le \emph{même} résultat. Il faut aussi la
\emph{confluence}, appelée propriété de
Church--Rosser : si $N \red N_1$ et $N \red N_2$,
il existe un terme~$N'$ tel que $N_1 \red N'$
et $N_2 \red N'$. Comme $N \red N'$ admet aussi
la suite vide, $\red$ est la clôture réflexive et
transitive de~$\redone$. Si $N_1$ et $N_2$ sont
normaux, leur seul réduit commun est eux-mêmes ;
ils sont donc égaux à renommage des variables liées
près.

\begin{prop}[Propriété locale de Church--Rosser]
  Pour les trois contractions principales et la
  conversion de permutation de la projection après
  analyse par cas explicitées dans la section précédente,
  si $N \redone N_1$ et $N \redone N_2$, alors
  il existe $N'$ tel que $N_1 \red N'$ et
  $N_2 \red N'$.
\end{prop}
\begin{proof}
  Deux redex situés dans des sous-termes disjoints
  se contractent dans l'un ou l'autre ordre et
  aboutissent au même terme, à renommage des variables
  liées près. Si le redex intérieur se trouve dans une
  partie recopiée par la contraction extérieure, on
  contracte ensuite \emph{toutes} ses copies résiduelles ;
  s'il se trouve dans une partie éliminée, aucune
  contraction résiduelle n'est nécessaire.
  Par exemple, dans $\proj{1}{\pair{O_1}{O_2}}$,
  une réduction à l'intérieur de~$O_1$ se retrouve
  après projection dans~$O_1'$ ; une réduction à
  l'intérieur de~$O_2$ est effacée par la projection,
  et les deux chemins atteignent~$O_1$.

  Il reste le chevauchement entre une analyse par cas
  d'une injection et la projection extérieure qui
  se permute avec cette analyse. Contracter d'abord
  l'analyse, puis projeter, donne la projection de
  la branche choisie après substitution. Permuter
  d'abord la projection dans les deux branches, puis
  contracter l'analyse, donne exactement ce même terme :
  la substitution sans capture commute avec cette
  projection. Les autres chevauchements des règles
  énoncées sont des cas de réduction dans une
  sous-expression ; l'argument des résidus les traite.
  Pour une famille plus large de permutations, une
  vérification distincte est nécessaire ; voir
  \href{https://jssst.or.jp/files/user/taikai/2015/PPL/ppl4-3.pdf}
  {Nakazawa et Fujita, \emph{Compositional Z:
  Confluence Proofs for Permutative Conversion},
  §\,3}.
\end{proof}

\begin{lem}[Lemme de Newman]
  Si $\redone$ est fortement normalisante et localement
  confluente, sa clôture réflexive et transitive~$\red$
  est confluente.
\end{lem}
\begin{proof}
  Procédons par récurrence bien fondée sur la relation
  $\redone$, possible grâce à la normalisation forte.
  Soient deux suites issues de~$N$ :
  \begin{align*}
    N &= N_0' \redone N_1' \redone \dots
      \redone N_k',\\
    N &= N_0'' \redone N_1'' \redone \dots
      \redone N_l''.
  \end{align*}
  Si l'une est vide, l'extrémité de l'autre
  est déjà un réduit commun.
  Sinon, leurs premiers pas
  donnent $N\redone N_1$ et $N\redone N_2$.
  Par confluence locale, $N_1$ et $N_2$ ont un réduit
  commun~$R$. L'hypothèse de récurrence appliquée
  à~$N_1$ et à~$N_2$ raccorde chaque suite initiale
  à un descendant de~$R$ ; appliquée ensuite à~$R$,
  elle raccorde ces deux descendants. Les extrémités
  des deux suites issues de~$N$ ont donc un réduit commun.
\end{proof}

Le théorème de normalisation forte et la propriété locale
précédente donnent ainsi la confluence pour les conversions
explicitement spécifiées. Chaque terme bien typé a une
forme normale unique à renommage des variables liées près.
Cette unicité concerne la relation de réduction fixée ;
ajouter d'autres conversions requiert de vérifier de
nouveau les hypothèses de confluence.

\end{document}
